\section{Distinctions Used in the Proof}
\label{app:history}

The following distinctions organize the arguments compared in
\cref{sec:experience}. Each row names a tempting simplification that fails,
its consequence for the proof, and where the proof handles it.
\begin{center}
\small
\begin{tabular}{@{}>{\raggedright\arraybackslash}p{.31\linewidth}>{\raggedright\arraybackslash}p{.40\linewidth}>{\raggedright\arraybackslash}p{.21\linewidth}@{}}
\toprule
Distinction & Consequence & Where\\
\midrule
Preserving the final reader state does not preserve the absence of
admitted comparisons. &
Compression shortens all blocks of a level simultaneously and records
admissions and block equalities as well as transitions. &
\cref{ex:loop-deletion}; \cref{sec:search}\\
\addlinespace
Betweenness does not make every continuation of a rejected word rejected. &
The tail proof works inside the graph of reachable rejected states. &
\cref{sec:tail}\\
\addlinespace
Shortening an arbitrary prefix of a forced tail can create an admitted pair
of equal suffixes. &
The proof replaces the prefix before the earliest periodic onset and keeps
the letter just before that onset. &
\cref{lem:onset,lem:no-infinite-prime}\\
\addlinespace
A finite solution need not refute the query. &
It serves only to exclude a cycle clash in the base system; the four-spine
extension supplies the finite countermodel. &
\cref{thm:finite-transfer}\\
\bottomrule
\end{tabular}
\end{center}
